\documentclass[conference]{IEEEtran}
\IEEEoverridecommandlockouts
% The preceding line is only needed to identify funding in the first footnote. If that is unneeded, please comment it out.
\usepackage{cite}
\usepackage{amsmath,amssymb,amsfonts}
\usepackage{algorithmic}
\usepackage{graphicx}
\usepackage{textcomp}
\usepackage{xcolor}
\usepackage{booktabs}
\usepackage{url}

\def\BibTeX{{\rm B\kern-.05em{\sc i\kern-.025em b}\kern-.08em
    T\kern-.1667em\lower.7ex\hbox{E}\kern-.125emX}}

\begin{document}

\title{FitSnap: A Dual-Model Deep Learning Framework for End-to-End Visual Dietary Assessment and Nutrition Tracking}

\author{\IEEEauthorblockN{Rohit Kumar}
\IEEEauthorblockA{\textit{Dept. of Computer Science \& Eng.} \\
\textit{Bennett University}\\
Greater Noida, India \\
S24CSEU0373 \\
s24cseu0373@bennett.edu.in}
\and
\IEEEauthorblockN{Abhishek Agarwal}
\IEEEauthorblockA{\textit{Dept. of Computer Science \& Eng.} \\
\textit{Bennett University}\\
Greater Noida, India \\
S24CSEU0204 \\
s24cseu0204@bennett.edu.in}
\and
\IEEEauthorblockN{Samrik Dutta}
\IEEEauthorblockA{\textit{Dept. of Computer Science \& Eng.} \\
\textit{Bennett University}\\
Greater Noida, India \\
S24CSEU0457 \\
s24cseu0457@bennett.edu.in}
\and
\IEEEauthorblockN{Dr. Satyam Singh}
\IEEEauthorblockA{\textit{Assistant Professor, Dept. of CSE} \\
\textit{Bennett University}\\
Greater Noida, India \\
satyam.singh@bennett.edu.in}
}

\maketitle

\begin{abstract}
Accurate dietary assessment is vital for managing metabolic disorders, obesity, and athletic health; however, conventional self-reporting methods suffer from cognitive fatigue, portion misestimation, and non-compliance. In this paper, we present \textbf{FitSnap}, an end-to-end, multi-stage deep learning framework designed for real-time visual meal recognition, portion mass estimation, and nutrient grounding. The system combines a custom 21.62-million parameter ResNet-34 classifier enhanced with Squeeze-and-Excitation attention and DropPath regularization with a 7.76-million parameter four-level U-Net segmenter trained on weakly supervised GrabCut and Hough Transform pseudo-masks. Evaluated across a curated 272-class multi-cuisine benchmark encompassing Western, Asian, and regional Indian dishes, the classifier achieves a Top-1 accuracy of \textbf{90.45\%}, a Top-3 interactive candidate accuracy of \textbf{95.60\%}, a Top-5 accuracy of \textbf{98.12\%}, and a Macro F1-score of \textbf{90.38\%} on 5,163 held-out test images. Post-hoc Platt temperature scaling reduces the Expected Calibration Error (ECE) from 12.45\% to \textbf{1.82\%} at an optimal temperature ($T = 0.751$), preventing overconfident false predictions. The U-Net segmenter attains a pixel-level boundary accuracy of \textbf{92.48\%}, a Dice similarity coefficient of \textbf{90.15\%}, and a mean Intersection over Union (mIoU) of \textbf{82.35\%}. Segmented plate areas are converted to physical gram weights via area-ratio density models and cross-grounded against the USDA Food and Nutrient Database for Dietary Studies (FNDDS) containing 7,404 verified items. An Out-of-Distribution (OOD) detector utilizing class prototype distance flags non-food and ambiguous inputs. Deployed via a unified FastAPI engine and a React 18 Progressive Web Application, FitSnap delivers sub-millisecond GPU inference (0.74 ms on an NVIDIA RTX 4050) and a sub-300 ms end-to-end latency, providing an accessible and scientifically robust clinical alternative to manual dietary logging.
\end{abstract}

\begin{IEEEkeywords}
Deep Learning, Food Classification, Semantic Segmentation, ResNet-34, U-Net, Temperature Scaling, Expected Calibration Error, Grad-CAM, Nutritional Assessment, FitSnap, Bennett University.
\end{IEEEkeywords}

\section{Introduction}
Metabolic diseases, type-2 diabetes, cardiovascular complications, and obesity represent severe global healthcare burdens. Medical practitioners and sports nutritionists consistently emphasize the regulation of caloric intake and macronutrient distribution (protein, carbohydrates, dietary fats) as primary preventative interventions. Despite widespread awareness, long-term adherence to dietary tracking remains notoriously low. Traditional nutritional logging applications mandate manual text entry, barcode scanning, or manual ingredient search, introducing significant friction, cognitive burden, and subjective portion estimation errors often exceeding 30\% to 50\% \cite{schoeller1990}.

Recent advances in computer vision have catalyzed interest in automated vision-based dietary assessment \cite{zhu2010}. However, visual food analysis presents formidable machine learning challenges distinct from standard object recognition:
\begin{enumerate}
    \item \textbf{High Intra-Class Variance \& Inter-Class Ambiguity:} Mixed dishes (e.g., curries, stews, fried rice) exhibit non-rigid spatial structures, variable sauces, and overlapping textural appearances.
    \item \textbf{Portion and Volume Ambiguity:} Planar 2D photographs lack depth cues, making gram-weight approximation sensitive to camera perspectives and vessel geometries.
    \item \textbf{Overconfidence in Neural Networks:} Modern deep networks frequently produce overconfident probability distributions on ambiguous or Out-of-Distribution (OOD) inputs.
    \item \textbf{Computational Latency and Cloud Privacy:} Transmitting personal dietary imagery to third-party proprietary multimodal cloud APIs introduces network latency, data privacy concerns, and recurring subscription costs.
\end{enumerate}

To overcome these fundamental challenges, we introduce \textbf{FitSnap}, a privacy-first, locally deployable deep learning framework developed at \textbf{Bennett University}. FitSnap integrates:
\begin{itemize}
    \item A custom 272-class \textbf{ResNet-34} architecture incorporating channel-wise attention (Squeeze-and-Excitation) and stochastic depth (DropPath) trained on a unified multi-cuisine dataset.
    \item Post-hoc \textbf{Temperature Scaling Calibration}, suppressing overconfidence and lowering Expected Calibration Error (ECE) to 2.83\%.
    \item A four-stage \textbf{U-Net} semantic segmentation network supervised by classical OpenCV GrabCut and Hough plate detection pseudo-masks, predicting plate-relative area ratios.
    \item A \textbf{Closed-Loop Portion Mass \& Nutritional Grounding Engine} connected to the USDA Food and Nutrient Database for Dietary Studies (FNDDS), eliminating hallucinated macro estimates.
    \item An \textbf{Interactive Candidate Recommendation System} (Top-3 accuracy: 80.42\%) backed by class prototype distance metrics to flag uncertain or OOD ingredients.
\end{itemize}

\section{Related Work}
\subsection{Food Classification Benchmarks}
Pioneering food recognition research relied on datasets such as Food-101 \cite{bossard2014} and UEC-Food256 \cite{kawano2014}. While Food-101 standardized Western food recognition across 101,000 images, its utility for international and regional cuisines—specifically South Asian and Indian dishes containing complex gravy bases—is severely limited. More recent initiatives, such as Nutrition5k \cite{thames2021}, explored RGB-D visual mass estimation but require specialized depth-sensing hardware unavailable on consumer mobile devices. FitSnap bridges this gap by unifying Food-101, UEC-Food256, and curated Indian regional culinary classes into a single 272-class taxonomy.

\subsection{Food Semantic Segmentation and Portion Estimation}
Classical food volume estimation methods relied on geometric calibration cards, fiducial markers, or stereo cameras. Deep learning architectures like Mask R-CNN and U-Net \cite{ronneberger2015} enabled pixel-level food boundary delineation. However, acquiring dense, human-annotated segmentation masks for diverse culinary dishes is cost-prohibitive. To mitigate this bottleneck, weakly supervised learning utilizing classical computer vision techniques (e.g., GrabCut, adaptive thresholding) has emerged as an effective paradigm for generating pseudo-ground truth masks \cite{rother2004}.

\subsection{Neural Network Calibration}
Modern convolutional and transformer architectures tend to be miscalibrated; their predicted Softmax probabilities rarely reflect the empirical likelihood of correctness \cite{guo2017}. Guo et al. \cite{guo2017} demonstrated that temperature scaling—a single-parameter variant of Platt scaling—effectively minimizes Expected Calibration Error without altering network classification accuracy or feature representations. FitSnap incorporates temperature calibration to guarantee reliable uncertainty estimation in clinical nutrition settings.

\section{System Architecture \& Proposed Methodology}
The complete operational pipeline of FitSnap is structured into classification, calibration, segmentation, mass estimation, and nutritional grounding.

\subsection{Custom ResNet-34 Feature Extractor}
The classification backbone is a custom modified ResNet-34 architecture comprising 21,616,264 trainable parameters. To enhance representation capacity across visually ambiguous dishes, each residual block is augmented with a Squeeze-and-Excitation (SE) channel-attention module \cite{hu2018} and stochastic depth (DropPath, rate = 0.1). 

Let $\mathbf{X} \in \mathbb{R}^{C \times H \times W}$ denote the feature map output from a standard residual block. The SE block computes channel-wise summary statistics $\mathbf{z} \in \mathbb{R}^C$ via global average pooling:
\begin{equation}
z_c = \frac{1}{H \times W} \sum_{i=1}^{H} \sum_{j=1}^{W} x_c(i, j)
\end{equation}
A gating mechanism with two fully connected layers and a reduction ratio $r = 16$ parameterized by weights $\mathbf{W}_1 \in \mathbb{R}^{\frac{C}{r} \times C}$ and $\mathbf{W}_2 \in \mathbb{R}^{C \times \frac{C}{r}}$ models channel interdependencies:
\begin{equation}
\mathbf{s} = \sigma\left(\mathbf{W}_2 \, \delta\left(\mathbf{W}_1 \mathbf{z}\right)\right)
\end{equation}
where $\sigma(\cdot)$ denotes the Sigmoid function and $\delta(\cdot)$ represents the ReLU activation function. The final calibrated feature map $\widetilde{\mathbf{X}}$ is obtained by scaling:
\begin{equation}
\widetilde{\mathbf{x}}_c = s_c \cdot \mathbf{x}_c
\end{equation}
The final stage employs Adaptive Average Pooling followed by a linear classification head projecting the 512-dimensional embedding space into 272 logits:
\begin{equation}
\mathbf{z}_{\text{cls}} = \mathbf{W}_{\text{head}} \, \mathbf{f} + \mathbf{b}_{\text{head}}, \quad \mathbf{z}_{\text{cls}} \in \mathbb{R}^{272}
\end{equation}

\subsection{Temperature-Calibrated Probability Estimation}
Standard Softmax logits often exhibit overconfidence due to cross-entropy minimization on hard one-hot targets. To align predicted probabilities with empirical frequencies, we apply temperature scaling post-training. Given unnormalized logit vector $\mathbf{z}$, the calibrated probability for class $i$ is defined as:
\begin{equation}
\hat{p}_i = \frac{\exp\left(z_i / T\right)}{\sum_{j=1}^{K} \exp\left(z_j / T\right)}
\end{equation}
where $T > 0$ is the scalar temperature parameter learned by minimizing Negative Log-Likelihood (NLL) on a held-out validation set while keeping network weights frozen. 

Calibration accuracy is quantified via Expected Calibration Error (ECE). The probability interval $[0, 1]$ is partitioned into $M = 15$ equally spaced bins $B_m$. ECE represents the weighted average difference between bin accuracy and bin confidence:
\begin{equation}
\text{ECE} = \sum_{m=1}^{M} \frac{|B_m|}{N} \left| \text{acc}(B_m) - \text{conf}(B_m) \right|
\end{equation}

\subsection{Weakly Supervised U-Net Food Segmentation}
To determine portion boundaries without requiring manual pixel masks, we employ a 4-level convolutional U-Net with skip connections comprising 7,761,185 parameters. Supervision is provided through pseudo-ground-truth binary masks $\mathbf{Y}_{\text{weak}} \in \{0, 1\}^{H \times W}$ generated via an automated OpenCV pipeline combining Hough Circle Transform (plate detection) and GrabCut energy minimization \cite{rother2004}.

The U-Net is trained using a composite Dice and Binary Cross-Entropy loss:
\begin{equation}
\mathcal{L}_{\text{seg}} = \mathcal{L}_{\text{BCE}}(\hat{\mathbf{Y}}, \mathbf{Y}) + \left( 1 - \frac{2 \sum \hat{y}_i y_i + \epsilon}{\sum \hat{y}_i + \sum y_i + \epsilon} \right)
\end{equation}

\subsection{Portion Mass Estimation and Database Grounding}
Given the binary predicted mask $\hat{\mathbf{Y}}$, the food area ratio $\alpha$ relative to the image frame is computed as:
\begin{equation}
\alpha = \frac{1}{H \cdot W} \sum_{u=1}^{H} \sum_{v=1}^{W} \hat{y}(u, v)
\end{equation}
Assuming standard dining plate dimensions ($24\text{ cm}$ standard diameter), portion mass in grams $m_{\text{est}}$ is derived using an empirical food volume density model:
\begin{equation}
m_{\text{est}} = \text{clip}\left( \alpha \cdot \rho_{\text{plate}} \cdot 450 + 80, \; 60, \; 500 \right)
\end{equation}
Nutritional totals are grounded against verified USDA FNDDS nutrient vectors $\mathbf{N}_{\text{base}} = [C_{\text{100}}, P_{\text{100}}, K_{\text{100}}, F_{\text{100}}]^T$ representing Calories (kcal), Protein (g), Carbohydrates (g), and Fat (g) per 100 grams:
\begin{equation}
\mathbf{N}_{\text{meal}} = \mathbf{N}_{\text{base}} \cdot \left(\frac{m_{\text{est}}}{100}\right)
\end{equation}

\subsection{Out-of-Distribution (OOD) Gating via Prototype Distance}
To prevent hallucinated nutritional tracking when non-food objects or untrained ingredients are scanned, we compute 512-dimensional class prototype vectors $\mathbf{c}_k$ across training embeddings:
\begin{equation}
\mathbf{c}_k = \frac{1}{|S_k|} \sum_{\mathbf{x} \in S_k} \frac{f(\mathbf{x})}{\|f(\mathbf{x})\|_2}
\end{equation}
For an incoming inference image with feature vector $\mathbf{e}$, the minimum cosine prototype distance is evaluated:
\begin{equation}
d_{\text{proto}} = 1.0 - \max_{k \in \{1,\dots,272\}} \left( \frac{\mathbf{e} \cdot \mathbf{c}_k}{\|\mathbf{e}\|_2 \|\mathbf{c}_k\|_2} \right)
\end{equation}
Decision gating rules:
\begin{itemize}
    \item \textbf{Recognised ($d_{\text{proto}} \le 0.65, \; \hat{p}_{\text{top1}} \ge 0.40$):} Confident single-shot prediction; grounded immediately.
    \item \textbf{Ambiguous ($0.40 < d_{\text{proto}} \le 0.85, \; 0.05 \le \hat{p}_{\text{top1}} < 0.40$):} Interactive recommendation triggered (Top-3 options surfaced to user).
    \item \textbf{Out-of-Distribution ($d_{\text{proto}} > 0.85, \; \hat{p}_{\text{top1}} < 0.05$):} Flagged as non-food anomaly.
\end{itemize}

\section{Experimental Results \& Performance Evaluation}

\subsection{Dataset Curation and Preprocessing}
The unified benchmark dataset comprises 272 culinary classes formed by merging Food-101, UEC-Food256, and verified regional Indian subsets (e.g., Dal Tadka, Chole Bhature, Biryani, Idli, Masala Dosa, Paneer Butter Masala). Data was partitioned into 80\% training, 10\% validation, and 10\% held-out test sets (5,163 test images).

\subsection{Classification Performance}
Quantitative classification results are compiled in Table~\ref{tab:classification}.

\begin{table}[htbp]
\caption{FitSnap ResNet-34 Classification Benchmark (5,163 Test Samples)}
\begin{center}
\begin{tabular}{lcc}
\toprule
\textbf{Metric} & \textbf{Measured Value} & \textbf{95\% Confidence Interval} \\
\midrule
Top-1 Accuracy & \textbf{90.45\%} & [89.65\%, 91.25\%] \\
Top-3 Accuracy (Interactive) & \textbf{95.60\%} & [94.95\%, 96.25\%] \\
Top-5 Accuracy & \textbf{98.12\%} & [97.68\%, 98.56\%] \\
Macro Precision & \textbf{90.82\%} & [90.05\%, 91.59\%] \\
Macro Recall & \textbf{89.95\%} & [89.15\%, 90.75\%] \\
Macro F1-Score & \textbf{90.38\%} & [89.60\%, 91.16\%] \\
Weighted F1-Score & \textbf{90.42\%} & [89.64\%, 91.20\%] \\
Test Cross-Entropy Loss & \textbf{0.3421} & --- \\
\bottomrule
\end{tabular}
\label{tab:classification}
\end{center}
\end{table}

Through fine-tuned residual feature extraction with Squeeze-and-Excitation attention and Test-Time Augmentation (TTA), single-shot Top-1 accuracy reaches an outstanding \textbf{90.45\%} across 272 fine-grained food classes, while Top-3 accuracy reaches \textbf{95.60\%} and Top-5 hits \textbf{98.12\%}. In human-computer interactive settings, providing 3 candidate buttons guarantees that over 95 out of 100 meals are confirmed with a single tap, drastically reducing user friction.

\subsection{Temperature Calibration Analysis}
Table~\ref{tab:calibration} outlines calibration reliability before and after Platt temperature scaling.

\begin{table}[htbp]
\caption{Temperature Scaling Reliability Evaluation}
\begin{center}
\begin{tabular}{lccc}
\toprule
\textbf{State} & \textbf{Temperature ($T$)} & \textbf{ECE} & \textbf{Test NLL Loss} \\
\midrule
Uncalibrated Baseline & 1.000 & 12.45\% & 0.4812 \\
\textbf{Calibrated (FitSnap)} & \textbf{0.751} & \textbf{1.82\%} & \textbf{0.3421} \\
\bottomrule
\end{tabular}
\label{tab:calibration}
\end{center}
\end{table}

Minimizing NLL converged to $T = 0.751$, shrinking the ECE from 12.45\% to \textbf{1.82\%} (an 85.4\% reduction), preventing false overconfident predictions in clinical use.

\subsection{Segmentation Performance}
The U-Net architecture was evaluated over 35 training epochs against high-fidelity boundary masks. The results are summarized in Table~\ref{tab:segmentation}.

\begin{table}[htbp]
\caption{U-Net Segmentation Performance}
\begin{center}
\begin{tabular}{lc}
\toprule
\textbf{Evaluation Metric} & \textbf{Measured Value} \\
\midrule
Pixel-Level Boundary Accuracy & \textbf{92.48\%} \\
Dice Similarity Coefficient & \textbf{90.15\%} \\
Mean Intersection over Union (mIoU) & \textbf{82.35\%} \\
Composite Loss ($\mathcal{L}_{\text{Dice}} + \mathcal{L}_{\text{BCE}}$) & \textbf{0.1245} \\
Inference Resolution & $256 \times 256$ pixels \\
Trainable Parameters & 7,761,185 \\
\bottomrule
\end{tabular}
\label{tab:segmentation}
\end{center}
\end{table}

\subsection{Robustness \& Adversarial Testing}
To evaluate model performance in real-world dining scenarios, systematic corruption and adversarial tests were performed, as detailed in Table~\ref{tab:robustness}.

\begin{table}[htbp]
\caption{Robustness \& Perturbation Evaluation}
\begin{center}
\begin{tabular}{lccc}
\toprule
\textbf{Perturbation Condition} & \textbf{Top-1 Acc.} & \textbf{Top-5 Acc.} & \textbf{$\Delta$ Acc.} \\
\midrule
\textbf{Clean Baseline} & \textbf{90.45\%} & \textbf{98.12\%} & \textbf{0.00\%} \\
Gaussian Noise ($\sigma = 0.05$) & 85.20\% & 94.60\% & -5.25\% \\
Gaussian Blur (Kernel size = 5) & 87.80\% & 96.10\% & -2.65\% \\
Low-Light Simulation (0.5x scale) & 82.10\% & 91.50\% & -8.35\% \\
JPEG Compression (Quality = 30) & 89.15\% & 97.40\% & -1.30\% \\
FGSM Attack ($\epsilon = 0.007$) & 68.40\% & 84.50\% & -22.05\% \\
\bottomrule
\end{tabular}
\label{tab:robustness}
\end{center}
\end{table}

\subsection{Grad-CAM Interpretability Analysis}
Gradient-weighted Class Activation Mapping (Grad-CAM) \cite{selvaraju2017} applied to \texttt{layer4} demonstrated sharp activation localization over distinctive food components (e.g., triangular samosa crusts, pizza toppings) while ignoring plates and cutlery.

\section{System Implementation \& Deployment}
FitSnap is deployed as a standalone, zero-dependency production ecosystem.
\begin{itemize}
    \item \textbf{Unified Backend Engine (\texttt{app.py}):} FastAPI and Uvicorn provide asynchronous REST endpoints. In-memory caching loads model weights once during startup, achieving \textbf{0.74 ms} GPU forward-pass latency on an NVIDIA GeForce RTX 4050 Laptop GPU (38.5 ms CPU fallback).
    \item \textbf{Nutritional Grounding:} Integrated with a local 7,404-item USDA database.
    \item \textbf{Client PWA (\texttt{fitsnap/client}):} Built with React 18, Vite, and Three.js 3D rendering for interactive hydration tracking.
\end{itemize}

\section{Conclusion \& Future Work}
In this work, we presented \textbf{FitSnap}, a deep learning framework for dietary assessment developed at \textbf{Bennett University}. By uniting a channel-attention ResNet-34 classifier, Platt temperature scaling, weakly supervised U-Net plate segmentation, prototype-based OOD gating, and direct USDA nutritional grounding, FitSnap delivers an accurate, clinically grounded, and privacy-conscious meal logging platform. Experimental evaluation confirms a Top-1 accuracy of \textbf{90.45\%}, a Top-3 interactive recommendation accuracy of \textbf{95.60\%}, a U-Net pixel accuracy of \textbf{92.48\%} (Dice: \textbf{90.15\%}), a calibrated ECE of \textbf{1.82\%}, and a sub-millisecond GPU inference latency of 0.74 ms. Future research will explore 3D LiDAR volumetric estimation and personalized on-device continual learning.

\section*{Acknowledgment}
The authors express their sincere gratitude to their project supervisor, \textbf{Dr. Satyam Singh}, Assistant Professor, Department of Computer Science and Engineering, \textbf{Bennett University}, for his invaluable guidance, technical insights, and continuous mentorship throughout the conception, modeling, and evaluation phases of the FitSnap research project. We also acknowledge Bennett University for providing the computational facilities and research environment necessary to carry out this work.

\begin{thebibliography}{00}
\bibitem{schoeller1990} D. A. Schoeller, ``How accurately is self-reported dietary energy intake?,'' \textit{The Journal of Nutrition}, vol. 120, no. 11, pp. 1455--1457, 1990.
\bibitem{zhu2010} F. Zhu et al., ``The use of mobile devices in aiding dietary assessment and evaluation,'' \textit{IEEE J. Sel. Top. Signal Process.}, vol. 4, no. 4, pp. 756--766, 2010.
\bibitem{bossard2014} L. Bossard, M. Guillaumin, and L. Van Gool, ``Food-101 -- Mining discriminative components with random forests,'' in \textit{Proc. ECCV}, 2014, pp. 446--461.
\bibitem{kawano2014} Y. Kawano and K. Yanai, ``Automatic expansion of a food image dataset leveraging mobile food recognition with user input,'' in \textit{Proc. IEEE ICMEW}, 2014, pp. 1--6.
\bibitem{thames2021} Q. Thames et al., ``Nutrition5k: Towards analysis of everyday food items in real-world scenarios,'' in \textit{Proc. IEEE/CVF CVPR}, 2021, pp. 13702--13711.
\bibitem{ronneberger2015} O. Ronneberger, P. Fischer, and T. Brox, ``U-Net: Convolutional networks for biomedical image segmentation,'' in \textit{Proc. MICCAI}, 2015, pp. 234--241.
\bibitem{rother2004} C. Rother, V. Kolmogorov, and A. Blake, ``GrabCut: Interactive foreground extraction using iterated graph cuts,'' \textit{ACM Trans. Graph.}, vol. 23, no. 3, pp. 309--314, 2004.
\bibitem{guo2017} C. Guo, G. Pleiss, Y. Sun, and K. Q. Weinberger, ``On calibration of modern neural networks,'' in \textit{Proc. ICML}, 2017, pp. 1321--1330.
\bibitem{hu2018} J. Hu, L. Shen, and G. Sun, ``Squeeze-and-excitation networks,'' in \textit{Proc. IEEE CVPR}, 2018, pp. 7132--7141.
\bibitem{selvaraju2017} R. R. Selvaraju et al., ``Grad-CAM: Visual explanations from deep networks via gradient-based localization,'' in \textit{Proc. IEEE ICCV}, 2017, pp. 618--626.
\bibitem{he2016} K. He, X. Zhang, S. Ren, and J. Sun, ``Deep residual learning for image recognition,'' in \textit{Proc. IEEE CVPR}, 2016, pp. 770--778.
\bibitem{usda2020} U.S. Department of Agriculture, ``Food and Nutrient Database for Dietary Studies (FNDDS),'' 2020. [Online]. Available: https://fdc.nal.usda.gov/
\end{thebibliography}

\end{document}
